\chapter{If and Match}

\textbf{What this chapter gives you:} the first tools that let a program \emph{decide}. You will write
comments, split long lines, branch with \code{if}/\code{elif}/\code{else}, test membership with \code{in},
and pick between patterns with \code{match}. Truthiness, \code{None}, and the one-line \code{a if c else b}
get extra attention, because they are the usual sources of silent bugs.

%RELATED: 2, 3, 4, 7, 9

% ------------------------------------------------------------
\section{Comment with \#}

A comment starts at \code{\#} and runs to the end of the line. Python ignores it. Use comments to say
\emph{why}, not to repeat what the code already says. There is no block comment: start every line with
\code{\#}. Inside a string, \code{\#} is just a character.

\begin{pycode}
>>> max_context = 200_000   # tokens the model can read at once
>>> note = "use # for comments"
>>> note
'use # for comments'
\end{pycode}

A triple-quoted string that stands alone is not a comment. It is a string value that Python builds and
throws away. Its proper use is as documentation (Chapter 14).

% ------------------------------------------------------------
\section{Continue lines with \textbackslash{} and brackets}

Keep lines under about 80 characters. If an expression is too long, you have two choices.

\begin{center}
\begin{tabularx}{\linewidth}{@{}l X l@{}}
\toprule
\textbf{Method} & \textbf{How} & \textbf{Verdict} \\
\midrule
Backslash & put \code{\textbackslash} as the very last character of the line & works, but fragile \\
Brackets  & break the line anywhere inside \code{( )}, \code{[ ]}, or \code{\{ \}} & preferred \\
\bottomrule
\end{tabularx}
\end{center}

\begin{pycode}
>>> total = 1200 + \
...         350 + \
...         40
>>> total
1590
>>> total = (1200 +
...          350 +
...          40)
>>> total
1590
\end{pycode}

\begin{trap}[: a space after the backslash]
A backslash followed by even one trailing space is a \code{SyntaxError}. You cannot see the space, which
makes it painful. Prefer brackets and you never meet this bug.
\end{trap}

Brackets also help with long prompts. Two string literals side by side are joined into one:

\begin{pycode}
>>> prompt = ("You are a helpful assistant. "
...           "Answer in one sentence.")
>>> prompt
'You are a helpful assistant. Answer in one sentence.'
\end{pycode}

\begin{trap}[: a missing comma joins strings]
In a list, \code{["a" "b"]} is a one-item list containing \code{"ab"}. A forgotten comma between two
string items silently merges them. You get no error, only a shorter list.
\end{trap}

% ------------------------------------------------------------
\section{Compare with if, elif, and else}

\subsection{The shape of an if statement}

An \code{if} line ends with a colon. The lines that belong to it are indented by the same amount, and
the standard is four spaces. Never mix tabs and spaces. Indentation \emph{is} the structure: when the
indent goes back, the block ends.

%MD<
% ```mermaid
% flowchart TD
%     A{"status == 200?"} -- True --> A1["print('ok')"]
%     A -- False --> B{"status == 429?"}
%     B -- True --> B1["print('rate limited')"]
%     B -- False --> C{"status >= 500?"}
%     C -- True --> C1["print('server error')"]
%     C -- False --> D["else: print('other')"]
% ```
%MD>
\begin{center}
\begin{tikzpicture}[dec/.style={draw=cblue,minimum width=30mm,minimum height=6.5mm,fill=cblue!8,font=\ttfamily\scriptsize,inner sep=1pt},
  act/.style={draw=cgreen,fill=cgreen!10,rounded corners=2pt,font=\ttfamily\scriptsize,minimum height=6mm,inner sep=3pt},
  lab/.style={font=\sffamily\tiny,text=cgray},>=Stealth]
  \node[dec] (d1) at (0,0)    {status == 200};
  \node[dec] (d2) at (0,-1.5) {status == 429};
  \node[dec] (d3) at (0,-3.0) {status >= 500};
  \node[act] (e)  at (0,-4.4) {else: print("other")};
  \node[act] (a1) at (4.6,0)    {print("ok")};
  \node[act] (a2) at (4.6,-1.5) {print("rate limited")};
  \node[act] (a3) at (4.6,-3.0) {print("server error")};
  \draw[->] (d1) -- node[lab,right]{False} (d2);
  \draw[->] (d2) -- node[lab,right]{False} (d3);
  \draw[->] (d3) -- node[lab,right]{False} (e);
  \draw[->] (d1) -- node[lab,above]{True} (a1);
  \draw[->] (d2) -- node[lab,above]{True} (a2);
  \draw[->] (d3) -- node[lab,above]{True} (a3);
  \node[lab,anchor=west,align=left] at (6.8,-1.5) {tests run top to bottom;\\the first True wins;\\only one branch runs};
\end{tikzpicture}
\end{center}

\begin{pycode}
status = 429
if status == 200:
    print("ok")
elif status == 429:
    print("rate limited: wait and retry")
elif status >= 500:
    print("server error")
else:
    print("other:", status)
\end{pycode}
\begin{rawcode}
rate limited: wait and retry
\end{rawcode}

\code{elif} and \code{else} are optional. Tests inside tests are allowed, one more indent level each,
but flat \code{elif} chains read better than deep nesting.

\subsection{Comparison operators}

Every comparison produces a \code{bool}.

\begin{center}
\begin{tabularx}{\linewidth}{@{}l l X@{}}
\toprule
\textbf{Operator} & \textbf{Meaning} & \textbf{Watch for} \\
\midrule
\code{==} & equal in value        & \code{5 == 5.0} is \code{True}; \code{"5" == 5} is \code{False} \\
\code{!=} & not equal             & the opposite of \code{==} \\
\code{<} \code{<=} \code{>} \code{>=} & ordering & mixing \code{int} and \code{str} raises \code{TypeError} \\
\code{is} / \code{is not} & same object   & use it for \code{None}, not for numbers or strings \\
\code{=} & \emph{assignment}     & \code{if x = 5:} is a \code{SyntaxError} \\
\bottomrule
\end{tabularx}
\end{center}

\begin{pycode}
>>> 5 == 5.0, "5" == 5
(True, False)
>>> 3 < "4"
Traceback (most recent call last):
  ...
TypeError: '<' not supported between instances of 'int' and 'str'
\end{pycode}

\subsection{and, or, not}

Use \code{and}, \code{or}, and \code{not} to combine tests. Comparisons are worked out first, so no
parentheses are needed, but add them when a reader might hesitate. Order of strength: \code{not}, then
\code{and}, then \code{or}.

\begin{pycode}
>>> tokens, has_key = 800, True
>>> tokens < 1000 and has_key
True
>>> tokens > 1000 or not has_key
False
\end{pycode}

\code{and} and \code{or} are \emph{short-circuit} operators. They stop as soon as the answer is known, and
they return one of their two \emph{operands}, not necessarily a \code{bool}.

\begin{typecheck}[: what and / or / not return]
\begin{center}
\begin{tabularx}{\linewidth}{@{}l X l@{}}
\toprule
\textbf{Expression} & \textbf{Rule} & \textbf{Type back} \\
\midrule
\code{a and b} & if \code{a} is falsy, return \code{a}; otherwise return \code{b} & type of an operand \\
\code{a or b}  & if \code{a} is truthy, return \code{a}; otherwise return \code{b} & type of an operand \\
\code{not a}   & \code{True} if \code{a} is falsy, else \code{False} & always \code{bool} \\
\bottomrule
\end{tabularx}
\end{center}
\end{typecheck}

\begin{pycode}
>>> "" or "default"
'default'
>>> "hello" or "default"
'hello'
>>> "hi" and 42
42
>>> 0 and 1 / 0
0
>>> not []
True
\end{pycode}

The \code{0 and 1 / 0} line never divides, because the left side already decided the result.

\begin{trap}[: \texttt{x == "a" or "b"} is always truthy]
\code{model == "claude" or "gemini"} means \code{(model == "claude") or "gemini"}. The right side is a
non-empty string, so the whole expression is never falsy. Each side of \code{or} must be a full test, or
use \code{in}.
\end{trap}

\begin{pycode}
>>> model = "gpt"
>>> model == "claude" or "gemini"
'gemini'
>>> model == "claude" or model == "gemini"
False
>>> model in ("claude", "gemini")
False
\end{pycode}

\subsection{Chained comparisons}

Python lets you write a range test the way you would on paper. \code{lo < x < hi} means
\code{lo < x and x < hi}, and \code{x} is evaluated once.

%MD<
% ```text
%  4    5         7                    10   11
%  |----o=========X====================o----|     5 < x < 10  with x = 7
%       (5 and 10 are excluded: open ends)
% ```
%MD>
\begin{center}
\begin{tikzpicture}[>=Stealth,font=\sffamily\scriptsize]
  \draw[->,cgray] (0,0) -- (8.2,0);
  \foreach \v/\x in {4/0,5/1.1,7/3.3,10/6.6,11/7.7} {
    \draw[cgray] (\x,-0.08) -- (\x,0.08);
    \node[below=2pt,font=\ttfamily\scriptsize] at (\x,-0.08) {\v};
  }
  \draw[line width=3pt,cgreen] (1.1,0) -- (6.6,0);
  \draw[fill=white,draw=cgreen,thick] (1.1,0) circle (0.09);
  \draw[fill=white,draw=cgreen,thick] (6.6,0) circle (0.09);
  \fill[cred] (3.3,0) circle (0.11);
  \node[above=4pt,text=cred] at (3.3,0.05) {x = 7};
  \node[above=4pt,text=cgreen!60!black] at (5.5,0.05) {\code{5 < x < 10} is True here};
\end{tikzpicture}
\end{center}

\begin{pycode}
>>> temperature = 0.7
>>> 0 <= temperature <= 2
True
>>> 5 < 7 < 10 < 999
True
\end{pycode}

\subsection{What is true?}

An \code{if} does not need a real \code{bool}. Python asks every object ``are you empty or zero?''.
Empty and zero things are \emph{falsy}; everything else is \emph{truthy}. \code{bool(x)} shows the verdict.

\begin{center}
\begin{tabularx}{\linewidth}{@{}l X@{}}
\toprule
\textbf{Falsy (counts as False)} & \textbf{Truthy (counts as True)} \\
\midrule
\code{False}, \code{None} & \code{True} \\
\code{0}, \code{0.0} & any other number, including \code{-1} and \code{0.1} \\
\code{""} (empty string) & any non-empty string, even \code{" "}, \code{"0"}, \code{"False"} \\
\code{[]}, \code{()}, \code{\{\}}, \code{set()}, \code{b""} & any container with at least one item, even \code{[0]} or \code{[None]} \\
\bottomrule
\end{tabularx}
\end{center}

\begin{pycode}
>>> [bool(v) for v in (0, 0.0, "", [], (), {}, set(), None)]
[False, False, False, False, False, False, False, False]
>>> [bool(v) for v in ("0", "False", " ", [0], [None], 0.1)]
[True, True, True, True, True, True]
\end{pycode}

\begin{gamerule}
Test emptiness the Python way: \code{if messages:} rather than \code{if len(messages) > 0:}. It also
works for strings, dicts, and sets.
\end{gamerule}

\subsection{None is not the same as empty or zero}

\code{None} means ``no value was given''. It is falsy, but so are \code{0} and \code{""}, and those
are real, valid values. To ask ``was anything given?'' test identity: \code{x is None} or
\code{x is not None}.

%MD<
% ```text
%            value        if value:      if value is not None:
%            None            False            False
%            0               False            True    <- valid temperature!
%            ""              False            True    <- valid empty string
%            0.7             True             True
% ```
%MD>
\begin{center}
\begin{tabularx}{0.8\linewidth}{@{}l c c@{}}
\toprule
\textbf{value} & \code{if value:} & \code{if value is not None:} \\
\midrule
\code{None} & \xmark\ False & \xmark\ False \\
\code{0}    & \xmark\ False & \cmark\ True (valid temperature) \\
\code{""}   & \xmark\ False & \cmark\ True (valid empty string) \\
\code{0.7}  & \cmark\ True  & \cmark\ True \\
\bottomrule
\end{tabularx}
\end{center}

\begin{boss}[: a falsy value that is perfectly valid]
A temperature of \code{0} is a real setting, yet it is falsy, so \code{if temperature:} treats it like
``not provided''. The same bug hides in \code{temperature or 1.0}, which swaps a valid \code{0} for
\code{1.0}.\\[2pt]
\textbf{Strategy:} when \code{None} means ``missing'', test \code{is None} / \code{is not None}. Use plain
truthiness only when ``empty or zero'' really should be treated as ``nothing''.
\end{boss}

\begin{pycode}
for temperature in (0, 0.7, None):
    wrong = temperature or 1.0
    right = 1.0 if temperature is None else temperature
    print(temperature, "wrong:", wrong, "right:", right)
\end{pycode}
\begin{rawcode}
0 wrong: 1.0 right: 0
0.7 wrong: 0.7 right: 0.7
None wrong: 1.0 right: 1.0
\end{rawcode}

(The \code{for} loop just repeats the block once per value; Chapter 7 explains it.)

\subsection{Conditional expressions: \texttt{a if cond else b}}

When you only need to \emph{choose a value}, a one-line form avoids a four-line \code{if}. The order is
unusual: the value for true comes first, the test is in the middle, and the value for false is last.
Python runs the test first.

%MD<
% ```text
%     "long"      if   tokens > 1000   else   "short"
%       ^                   ^                    ^
%   (2) used if         (1) test runs        (2) used if
%       True                first                False
% ```
%MD>
\begin{center}
\begin{tikzpicture}[blk/.style={draw=#1,fill=#1!10,rounded corners=2pt,font=\ttfamily\small,minimum height=7mm,inner sep=4pt},
  kw/.style={font=\ttfamily\small,text=cpurple},
  num/.style={circle,fill=cgold,text=white,font=\sffamily\scriptsize\bfseries,inner sep=1.5pt},
  lab/.style={font=\sffamily\scriptsize,text=cgray,align=center},>=Stealth]
  \node[blk=cgreen] (t) at (0,0) {"long"};
  \node[kw] (if) at (1.7,0) {if};
  \node[blk=cblue] (c) at (4.0,0) {tokens > 1000};
  \node[kw] (el) at (6.3,0) {else};
  \node[blk=cred] (f) at (7.9,0) {"short"};
  \node[num] at (4.0,1.0) {1};
  \node[lab] at (4.0,1.6) {test runs first};
  \node[num] at (0,1.0) {2};
  \node[lab] at (0,1.6) {result if True};
  \node[num] at (7.9,1.0) {2};
  \node[lab] at (7.9,1.6) {result if False};
  \draw[->,cgray] (c.north) -- (4.0,0.8);
\end{tikzpicture}
\end{center}

\begin{pycode}
>>> tokens = 1500
>>> label = "long" if tokens > 1000 else "short"
>>> label
'long'
>>> print("over limit" if tokens > 4096 else "fits")
fits
\end{pycode}

\begin{typecheck}[: what a conditional expression gives back]
It is an \emph{expression}: it produces one value, so you can assign it, pass it to a function, or put
it inside an f-string. The two branches may be different types; the result has the type of whichever
branch ran. It has \code{else} always, because it must produce a value.
\end{typecheck}

\begin{trap}[: written backwards, it runs and is wrong]
Putting the test first looks natural but is valid Python with a different meaning. The line below does not
raise an error. It treats \code{"long"} as the test, which is always truthy, so you get
\code{tokens > 1000} (a \code{bool}) instead of a label.
\end{trap}

\begin{pycode}
>>> tokens = 1500
>>> tokens > 1000 if "long" else "short"
True
>>> label = "long" if tokens > 1000 else "short"
\end{pycode}

\begin{trap}[: it binds loosely]
The whole thing is chosen \emph{after} the surrounding \code{+} is evaluated. Add parentheses when you mix
them.
\end{trap}

\begin{pycode}
>>> flag = False
>>> "Result: " + "yes" if flag else "no"
'no'
>>> "Result: " + ("yes" if flag else "no")
'Result: no'
\end{pycode}

\begin{gamerule}
Use \code{a if c else b} to pick between two values. Use a normal \code{if} block when you want to
\emph{do} something (call functions, change variables). Do not chain several in one line.
\end{gamerule}

% ------------------------------------------------------------
\section{Do multiple comparisons with in}

Many \code{or}-ed equality tests collapse into one \code{in}. \code{x in container} gives a \code{bool}.
What ``in'' means depends on the container.

\begin{center}
\begin{tabularx}{\linewidth}{@{}l X l@{}}
\toprule
\textbf{Container} & \textbf{\code{in} checks} & \textbf{Speed} \\
\midrule
\code{str}            & is it a \emph{substring}? (\code{"at" in "cat"} is True) & scans \\
\code{list}, \code{tuple} & is it equal to an item? (\code{"at" in ["cat"]} is False) & scans \\
\code{set}            & is it an item?   & instant \\
\code{dict}           & is it a \emph{key}? (values are ignored) & instant \\
\bottomrule
\end{tabularx}
\end{center}

\begin{pycode}
>>> model = "claude-haiku-5-5"
>>> "haiku" in model
True
>>> allowed = {"claude-haiku-5-5", "claude-opus-5-5"}
>>> model in allowed
True
>>> limits = {"input": 100, "output": 20}
>>> "input" in limits, 100 in limits
(True, False)
>>> "at" in "cat", "at" in ["cat"]
(True, False)
>>> "gpt" not in allowed
True
\end{pycode}

\begin{gamerule}
If you test membership many times against a fixed collection, store it in a \code{set}. Sets check
instantly; lists check one item at a time (Chapter 9).
\end{gamerule}

% ------------------------------------------------------------
\section{The walrus operator}

\code{name := expression} (Python 3.8+) assigns \emph{and} gives back the value, so you can test and
keep a result in one step. It looks like a walrus lying on its side.

\begin{pycode}
>>> reply = "x" * 120
>>> limit = 100
>>> if (over := len(reply) - limit) > 0:
...     print("trim", over, "characters")
trim 20 characters
>>> over
20
\end{pycode}

The variable stays alive after the \code{if}. It also works in \code{while} loops (Chapter 7).

\begin{trap}[: parentheses decide what is assigned]
Without parentheses around the assignment, the walrus grabs the whole comparison, so you store a
\code{bool}, not the number.
\end{trap}

\begin{pycode}
>>> reply, limit = "x" * 120, 100
>>> if over := len(reply) - limit > 0:
...     print(over)
True
\end{pycode}

% ------------------------------------------------------------
\section{Match}

\code{match} (Python 3.10+) compares one \emph{subject} against a list of \emph{patterns}, top to bottom.
The first pattern that fits runs its block and the whole \code{match} ends. There is no fall-through and
no \code{break}. \code{case \_} is the default.

%MD<
% ```mermaid
% flowchart TD
%     S(["match status:"]) --> C1{"case 200?"}
%     C1 -- yes --> B1["msg = 'ok'"]
%     C1 -- no --> C2{"case 401 | 403?"}
%     C2 -- yes --> B2["msg = 'not allowed'"]
%     C2 -- no --> C3{"case 404?"}
%     C3 -- yes --> B3["msg = 'missing'"]
%     C3 -- no --> C4["case _  (always fits)"] --> B4["msg = 'unknown'"]
% ```
%MD>
\begin{center}
\begin{tikzpicture}[cs/.style={draw=cblue,fill=cblue!8,font=\ttfamily\scriptsize,minimum width=26mm,minimum height=6.5mm,inner sep=2pt},
  bd/.style={draw=cgreen,fill=cgreen!10,rounded corners=2pt,font=\ttfamily\scriptsize,minimum height=6.5mm,inner sep=3pt},
  lab/.style={font=\sffamily\tiny,text=cgray},>=Stealth]
  \node[draw=cgray,rounded corners=6pt,font=\ttfamily\scriptsize,inner sep=3pt] (s) at (0,0.9) {match status:};
  \node[cs] (c1) at (0,0)    {case 200};
  \node[cs] (c2) at (0,-1.0) {case 401 | 403};
  \node[cs] (c3) at (0,-2.0) {case 404};
  \node[cs,fill=cgold!15,draw=cgold] (c4) at (0,-3.0) {case \_};
  \node[bd] (b1) at (4.4,0)    {msg = "ok"};
  \node[bd] (b2) at (4.4,-1.0) {msg = "not allowed"};
  \node[bd] (b3) at (4.4,-2.0) {msg = "missing"};
  \node[bd] (b4) at (4.4,-3.0) {msg = "unknown"};
  \draw[->] (s) -- (c1);
  \foreach \a/\b in {c1/c2,c2/c3,c3/c4} \draw[->] (\a) -- node[lab,right]{no} (\b);
  \foreach \a/\b in {c1/b1,c2/b2,c3/b3,c4/b4} \draw[->] (\a) -- node[lab,above]{fits} (\b);
  \node[lab,anchor=west,align=left] at (6.9,-1.5) {first fit wins;\\then the match is over};
\end{tikzpicture}
\end{center}

\subsection{Simple matches}

\begin{pycode}
status = 404
match status:
    case 200:
        msg = "ok"
    case 401 | 403:
        msg = "not allowed"
    case 404:
        msg = "missing"
    case _:
        msg = "unknown"
print(msg)
\end{pycode}
\begin{rawcode}
missing
\end{rawcode}

\code{401 | 403} is an \emph{or-pattern}: either value fits. The same job can be done with an
\code{if}/\code{elif} chain or, for pure lookup, a dictionary. Pick by readability.

\begin{center}
\begin{tabularx}{\linewidth}{@{}l X X@{}}
\toprule
\textbf{Tool} & \textbf{Best when} & \textbf{Default case} \\
\midrule
\code{if}/\code{elif}/\code{else} & conditions are ranges or complex tests & \code{else} \\
\code{match}/\code{case}         & you compare against shapes or exact values & \code{case \_} \\
\code{dict.get(key, default)}    & you only map a value to another value & the second argument \\
\bottomrule
\end{tabularx}
\end{center}

\begin{pycode}
>>> messages = {200: "ok", 404: "missing"}
>>> messages.get(500, "unknown")
'unknown'
\end{pycode}

\subsection{Structural matches}

``Structural'' means the pattern describes the \emph{shape} of the subject and can pull values out of it.
A sequence pattern matches a list or tuple of the same length. A name in a pattern \emph{captures} that
part into a variable. \code{*rest} captures the remaining items as a list. A \code{case} may end with
\code{if condition}, called a \emph{guard}.

\begin{pycode}
for line in ("quit", "add gpt-x", "add claude-a 3", "list all now"):
    parts = line.split()
    match parts:
        case ["quit"]:
            print("bye")
        case ["add", name]:
            print("add one:", name)
        case ["add", name, *rest]:
            print("add", name, "with extras", rest)
        case [cmd, *_] if cmd == "list":
            print("listing")
        case _:
            print("unknown")
\end{pycode}
\begin{rawcode}
bye
add one: gpt-x
add claude-a with extras ['3']
listing
\end{rawcode}

\begin{typecheck}[: pattern cheat sheet]
\begin{center}
\begin{tabularx}{\linewidth}{@{}l X l@{}}
\toprule
\textbf{Pattern} & \textbf{Fits when the subject is} & \textbf{Binds} \\
\midrule
\code{200}, \code{"ok"}, \code{None} & equal to that literal (\code{None} by identity) & nothing \\
\code{401 | 403}            & any one of the alternatives & nothing \\
\code{x}                    & \emph{anything} & \code{x} \\
\code{\_}                   & anything & nothing \\
\code{[a, b]} or \code{a, b} & a list or tuple of exactly 2 (not a \code{str}) & \code{a}, \code{b} \\
\code{[a, *rest]}           & a sequence of 1 or more & \code{a}, \code{rest} (list) \\
\code{\{"k": v\}}            & a dict with key \code{"k"}; extra keys are fine & \code{v} \\
\code{int()}, \code{str()}  & an instance of that type & nothing \\
\code{case p if cond}       & \code{p} fits and \code{cond} is true & as \code{p} \\
\bottomrule
\end{tabularx}
\end{center}
\end{typecheck}

\subsection{Matching JSON-like data}

Model replies and tool calls are dicts. A mapping pattern checks keys and values and pulls out what you
need. Unlike sequences, a dict may have extra keys.

\begin{pycode}
events = [
    {"type": "text", "text": "Hello"},
    {"type": "tool_use", "name": "search", "input": {"q": "tokens"}},
    {"type": "ping"},
    42,
]
for event in events:
    match event:
        case {"type": "text", "text": body}:
            print("text:", body)
        case {"type": "tool_use", "name": name}:
            print("tool call:", name)
        case {"type": other}:
            print("ignored:", other)
        case int() | float():
            print("a number")
\end{pycode}
\begin{rawcode}
text: Hello
tool call: search
ignored: ping
a number
\end{rawcode}

If no case fits and there is no \code{case \_}, the \code{match} simply does nothing. It is not an
error.

\begin{trap}[: a bare name captures, it does not compare]
In a pattern, a plain name such as \code{LIMIT} is a \emph{capture}: it always fits and overwrites the
variable. To compare against a constant, use a dotted name (\code{Config.LIMIT}), a literal, or a guard.
\end{trap}

\begin{pycode}
>>> LIMIT = 100
>>> match 5:
...     case LIMIT:
...         print("matched; LIMIT is now", LIMIT)
matched; LIMIT is now 5
>>> match 5:
...     case n if n == 100:
...         print("equal to 100")
...     case _:
...         print("not 100")
not 100
\end{pycode}

\begin{trap}[: names are bound even when the guard fails]
A pattern may capture a name and then lose to its guard. The capture stays behind.
\end{trap}

\begin{pycode}
>>> match [3, 5]:
...     case x, y if y > 6:
...         print("big")
...     case _:
...         print("no match")
no match
>>> x, y
(3, 5)
\end{pycode}

\begin{practice}
\begin{enumerate}
\item Write an \code{if}/\code{elif}/\code{else} that prints \code{"too low"}, \code{"too high"}, or
\code{"just right"} when \code{guess} is compared with \code{secret}.
\item What does \code{"yes" and 0} evaluate to? What does \code{"" or None} evaluate to?
\item Rewrite \code{label = "big" if n > 10 else "small"} as a normal four-line \code{if}/\code{else}.
Then explain what \code{n > 10 if "big" else "small"} would return.
\item Why is \code{if temperature:} wrong when 0 is a legal setting? Write the correct test.
\item Write a \code{match} that handles the strings \code{"start"} and \code{"stop"}, plus a default.
\item What is the difference between \code{"at" in "cat"} and \code{"at" in ["cat"]}?
\end{enumerate}
\end{practice}

\begin{achievement}{Fork in the Road}
You can now:
\begin{itemize}
\item add comments and split long lines with brackets;
\item build \code{if}/\code{elif}/\code{else} ladders and chained comparisons;
\item predict truthiness, and test \code{None} with \code{is None};
\item write \code{a if cond else b} in the right order, and spot it written backwards;
\item use \code{in}, the walrus, and \code{match} with sequence and mapping patterns.
\end{itemize}
\end{achievement}
